\documentclass[10pt,a4paper]{amsart}
\usepackage[margin=19mm]{geometry}
\usepackage{amsmath,amssymb,mathtools}
\usepackage[scr=rsfs,cal=euler]{mathalfa}
\usepackage{xfrac}
\usepackage[unicode,hidelinks,bookmarks=false]{hyperref}
\newcommand{\ie}{i.e.\ }
\newcommand{\cf}{cf.\ }
\newcommand{\resp}{respectively\ }
\newcommand{\Subsection}{\subsection}
\providecommand{\nobreakdash}{\nobreak\hbox{-}}
\setlength{\emergencystretch}{2em}
\multlinegap0pt
\usepackage{amssymb}
\usepackage{mathtools}
\usepackage{tikz}
\usepackage[shortlabels]{enumitem}
\usepackage{csquotes}
\newtheorem{proposition}{Proposition}
\newcommand{\Ah}{\widehat A}
\newcommand{\Ahcwline}{\widehat A\text{-}\mathrm{cw}^{\mathrm{line}}}
\DeclareMathOperator{\Scal}{Scal}
\DeclareMathOperator{\comass}{comass}
\newcommand{\ip}[2]{\left\langle #1,#2\right\rangle}
\newcommand{\norm}[1]{\left\lVert #1\right\rVert}
\title{Stable systolic inequalities via mod n covering}
\author{Aditya Kumar}
\date{Source: arXiv:2604.26891v2 (CC BY 4.0)}
\begin{document}
\maketitle

\noindent\emph{Source and licence.} Stable systolic inequalities via mod n covering. Version arXiv:2604.26891v2 (CC BY 4.0), distributed by arXiv under the Creative Commons Attribution 4.0 International licence, http://creativecommons.org/licenses/by/4.0/, as recorded in the arXiv licence metadata for this exact version and shown on the abstract page https://arxiv.org/abs/2604.26891v2. Full licence notice: \texttt{licenses/arxiv-2604.26891-CC-BY-4.0.txt}.\par\smallskip

\noindent\textit{Editorial scope.} Counted unit: the article's proposition prop:refined-cowaist (source0.tex lines 714--725). The article's complete original proof of it is delivered with it (source0.tex lines 727--797, one \texttt{proof} environment of 71 lines). No mathematics is written, repaired or re-derived: the statement and the proof are the article's own lines, in the article's own notation, and no retained line is substituted or re-flowed.  No further source-local context is needed: every cross-reference of every retained line resolves inside the two delivered spans.  Nothing was omitted from any retained span, so the package carries no omission record.  No retained line contains a citation, so the wrapper carries no bibliography and no bibliography entry.  No retained line contains a figure environment, an image-inclusion command or drawing code, so no image is delivered and the package declares no asset dependency.  Editorial formatting only, in addition: an AMS \texttt{amsart} preamble that restates the article's own declarations and macros used by the retained lines, namely the environments \texttt{proposition} on separate counters, together with the standard packages those lines need.  The retained environments print in this document as Proposition 1 in order of appearance, whereas the article prints the same items with the numbering of its own full document.  The complete native source of the article as distributed in the arXiv e-print for this exact version is bundled unchanged as \texttt{sources/199431/source0.tex}.
\begin{proposition}\label{prop:refined-cowaist}
Let \((M^{2m},g)\) be a closed spin Riemannian manifold, and let \(\sigma=\min_M\Scal_g>0\). Then
\[
\Ahcwline(M,g)
\le
\frac{4m}{\sigma}.
\]
Equivalently, every \(\Ah\)-admissible Hermitian line bundle \(L\to M\) with unitary connection satisfies
\begin{equation}\label{eq:waistequiv}
\norm{R^L}_\infty \ge \frac{\sigma}{4m}.
\end{equation}
\end{proposition}

\begin{proof}
Let \(L\to M\) be an \(\Ah\)-admissible Hermitian line bundle with a unitary connection. Since the twisted Dirac index is nonzero, the corresponding twisted Dirac operator has a nonzero harmonic spinor \(\psi\).

The Lichnerowicz formula gives
\[
0
=
\int_M
\left(
|\nabla\psi|^2
+
\frac{\Scal_g}{4}|\psi|^2
+
\ip{\mathcal R^L\psi}{\psi}
\right)dV,
\]
where
\[
\mathcal R^L
=
\frac12
\sum_{i,j}c(e_i)c(e_j)R^L(e_i,e_j).
\]
We claim that \(\norm{\mathcal R^L}_{\operatorname{op}}\le m\norm{R^L}_\infty\). Fix a point $x\in M$. Since $L_x$ is one-dimensional and the connection is unitary, write $R^L=i\alpha$ with $\alpha$ a real two-form at $x$. Choose an orthonormal basis putting $\alpha$ in skew-normal form:
\begin{equation}\label{eq:normal-form}
\alpha
=
\lambda_1 e^1\wedge e^2
+
\lambda_2 e^3\wedge e^4
+
\cdots+
\lambda_n e^{2m-1}\wedge e^{2m}.
\end{equation}
For a two-form in this normal form, its comass is $\max_j |\lambda_j|$. With our convention for the line bundle curvature norm, this means\( \norm{R^L}_x=\norm{\alpha}_{\comass}=\max_j|\lambda_j|.\)
In the basis \eqref{eq:normal-form}, the twisting term becomes
\begin{equation*}
\mathcal R^L
=
\sum_{j=1}^m i\lambda_j c(e_{2j-1})c(e_{2j}).
\end{equation*}
Set $B_j=i c(e_{2j-1})c(e_{2j})$. The operators $B_j$ are self-adjoint involutions, and they commute because they involve disjoint Clifford generators. Hence they can be diagonalized simultaneously, with eigenvalues $\pm1$. Therefore the eigenvalues of $\mathcal R^L$ are of the form $\sum_j \pm \lambda_j$. Therefore we have as claimed,
\begin{equation}\label{eq:op-norm-bound}
\norm{\mathcal R^L}_{\operatorname{op}}
\le
\sum_{j=1}^m |\lambda_j|
\le
m\max_j|\lambda_j|
=
m\norm{R^L}_x.
\end{equation}
It follows that
\[
\ip{\mathcal R^L\psi}{\psi}
\ge
-m\norm{R^L}_\infty|\psi|^2.
\]
Using \(\Scal_g\ge\sigma\), we obtain
\[
0
\ge
\int_M
\left(
|\nabla \psi|^2
+
\left(\frac{\sigma}{4}-m\norm{R^L}_\infty\right)
|\psi|^2
\right)dV.
\]
If \(\norm{R^L}_\infty<\sigma/(4m)\), then the integrand is nonnegative and is positive wherever \(\psi\neq0\), a contradiction. Hence \(\norm{R^L}_\infty\ge\sigma/(4m)\). Taking reciprocals and then taking the supremum over \(\Ah\)-admissible line bundles gives the cowaist inequality.
\end{proof}

\end{document}
